% SPDX-License-Identifier: CC-BY-4.0
\documentclass[11pt,letterpaper]{article}

\usepackage[T1]{fontenc}
\usepackage{amsmath,amssymb}
\usepackage{tikz}
\usepackage{pgfplots}
\pgfplotsset{compat=1.18}
\usepackage{algorithm}
\usepackage{algpseudocode}
\usepackage{booktabs}
\usepackage{enumitem}
\usepackage[hidelinks]{hyperref}

\newcommand{\N}{\mathbb{N}}
\newcommand{\Z}{\mathbb{Z}}
\newcommand{\Q}{\mathbb{Q}}
\newcommand{\R}{\mathbb{R}}
\newcommand{\C}{\mathbb{C}}
\newcommand{\Map}[3]{#1 \colon #2 \to #3}
\newcommand{\Mod}[1]{#1\text{-}\mathrm{Mod}}
\newcommand{\GroupHom}{\mathrm{Hom}_{\mathbf{Grp}}}
\newcommand{\RingHom}{\mathrm{Hom}_{\mathbf{Ring}}}

\newtheorem{theorem}{Theorem}
\newtheorem{lemma}[theorem]{Lemma}
\newtheorem{definition}[theorem]{Definition}
\newtheorem{remark}[theorem]{Remark}
\newenvironment{proof}{\noindent\textit{Proof.}}{\hfill\(\square\)}

\title{Classical absorption at one target}
\author{ascendnoosphere}
\date{\today}



\begin{document}
\maketitle

\begin{abstract}
The raw maps with coefficients \(3\cdot2^a,2^a3^b,2\) have the same
positive universal capture obligation as the classical raw map.
Binary compression preserves the exact odd target \(3^b\).
A finite prefix enters the invariant ternary lattice, after at most \(b\)
compressed affine instructions, and one common entry witness determines
the entire subsequent tail. We retain both the original raw clock and
the separate classical variable clock. The designated cycle has
primitive original period \(a+3\). No classical convergence hypothesis
is proved or assumed in the equivalence theorem.
\end{abstract}

\section{Parameters and the exact target}

Throughout, \(a,b\in\N\), with zero allowed, and \(\sigma=3^b\).
For every signed integer state, define
\[
F_{a,b}(n)=
\begin{cases}
n/2,&2\mid n,\\
2^a(3n+\sigma),&2\nmid n.
\end{cases}
\]
Division has precedence, including at zero. Put \(G_b=F_{0,b}\) and
\(H=F_{0,0}\). Thus \(H\) is the classical raw map, with odd instruction
\(3n+1\), not the shortcut map with odd instruction \((3n+1)/2\).
Iterates are indexed by \(\N\), including time zero.

The proofs below use exact coprime scaling
\cite{AscendScaling}, binary block compression \cite{AscendCompressionI},
and full raw-time target recovery \cite{AscendCompressionII}.
The universal assertions always quantify over positive seeds.
Pointwise compression identities retain every signed seed.

% Oddness is the branch information needed for both scaling and target reflection.
\begin{lemma}[Ternary scale]\label{ScaleOdd}
The integer \(\sigma\) is positive, satisfies \(\gcd(\sigma,2)=1\), and is odd.
For every \(m\in\Z\),
\[
G_b(\sigma m)=\sigma H(m),
\qquad
F_{a,b}(\sigma m)=\sigma F_{a,0}(m).
\]
Consequently, both equalities hold for every iterate with the same raw time.
\end{lemma}

\begin{proof}
Positivity follows from \(3>0\). Every power of \(3\) is coprime to \(2\),
so \(\sigma\) is odd. Hence \(2\mid\sigma m\) if and only if \(2\mid m\).
On this legal division branch, \(\sigma m/2=\sigma(m/2)\).
On the other branch,
\[
3\sigma m+\sigma=\sigma(3m+1),
\qquad
2^a(3\sigma m+\sigma)=\sigma\,2^a(3m+1).
\]
This proves both one-step identities. Induction on the number of raw
steps proves their iterate versions: substitute the preceding iterate
as \(m\) into the one-step identity. The image \(\sigma\Z\) is therefore
invariant. Multiplication by \(\sigma\) is a bijection from \(\Z\) to
this invariant lattice, not to the whole original state space.
\end{proof}

\section{Binary blocks and full target reflection}

Let \(R_b(k,n)\) count the affine instructions among the first \(k\)
steps of \(G_b\) from \(n\), and put
\[
\tau_{a,b}(k,n)=k+aR_b(k,n).
\]
This count refers to actual compressed endpoint states.

% Every block is legal, including when its affine endpoint is zero.
\begin{lemma}[Exact block clock]\label{BlockClock}
For every \(n\in\Z\) and \(k\in\N\),
\[
F_{a,b}^{\,\tau_{a,b}(k,n)}(n)=G_b^k(n).
\]
At an even compressed state a block has length one.
At an odd compressed state \(x\) a block has length \(a+1\), and
\[
F_{a,b}^{\,\ell}(x)=2^{a+1-\ell}(3x+\sigma)
\qquad(1\le\ell\le a+1).
\]
Every strict interior state of this latter block is even.
\end{lemma}

\begin{proof}
For even \(x\), both maps use the same division \(x/2\).
For odd \(x\), the first raw state is \(2^a(3x+\sigma)\).
For \(1\le\ell<a+1\), the displayed exponent is positive, so division by
\(2\) is legal and removes one binary factor. Induction proves the formula
through its endpoint \(3x+\sigma=G_b(x)\). The argument also applies when
that endpoint is zero, since division is then still legal.
When \(a=0\), the block consists only of its endpoint and has no interior.

Concatenate these blocks along the actual \(G_b\) orbit.
Each compressed division contributes one original step and each affine
instruction contributes \(a+1\). Induction on \(k\) gives the endpoint
identity and the clock formula.
\end{proof}

% Strictly positive clock increments account for every original raw time.
\begin{theorem}[Exact odd-target reflection]\label{TargetReflection}
For every signed \(n\),
\[
\bigl(\exists r\in\N:F_{a,b}^r(n)=\sigma\bigr)
\quad\Longleftrightarrow\quad
\bigl(\exists k\in\N:G_b^k(n)=\sigma\bigr).
\]
\end{theorem}

\begin{proof}
A compressed hit expands to a raw hit by Lemma~\ref{BlockClock}.
Conversely, consecutive clock values differ by one or by \(a+1\),
and \(\tau_{a,b}(k,n)\ge k\). For each raw time \(r\), there is a largest
\(k\) with \(\tau_{a,b}(k,n)\le r\): the set is nonempty because it contains
zero and finite because \(k\le r\). Thus
\[
\tau_{a,b}(k,n)\le r<\tau_{a,b}(k+1,n).
\]
If \(r\) is not the left endpoint, it is a strict block interior.
Such a state is even by Lemma~\ref{BlockClock}, whereas \(\sigma\) is odd
by Lemma~\ref{ScaleOdd}. A hit must therefore occur at an endpoint, which
is the same compressed state. This proves reflection at every raw time,
not only at a selected subsequence.
\end{proof}

\section{Finite entry and one common tail}

% The induction preserves the inherited divisibility obligation at every child.
\begin{lemma}[Finite ternary entry]\label{FiniteEntry}
For every \(n>0\) there are \(t\in\N\) and \(m>0\) such that
\[
G_b^t(n)=\sigma m,
\qquad
A_{G_b}(t,n)\le b.
\]
Here \(A_{G_b}\) counts the affine instructions in that compressed raw
prefix. It does not count divisions.
\end{lemma}

\begin{proof}
More generally, suppose \(x>0\), \(r\in\N\), and \(\sigma\mid3^r x\).
We prove finite entry with at most \(r\) future affine instructions by
well-founded induction on the lexicographic pair \((r,x)\in\N^2\).

If \(\sigma\mid x\), use time zero. If not, then \(r>0\), because the
case \(r=0\) would already give \(\sigma\mid x\).
If \(x\) is even, write \(x=2y\). Positivity gives \(0<y<x\).
Since \(\sigma\) is coprime to \(2\), the inherited relation
\(\sigma\mid 2\cdot3^r y\) implies \(\sigma\mid3^r y\).
The legal division thus preserves \(r\), strictly decreases \(x\), and
uses no affine instruction. Apply the induction hypothesis at \((r,y)\).

If \(x\) is odd, its next state \(y=3x+\sigma\) is positive.
The two relations
\[
\sigma\mid3^r x,\qquad
\sigma\mid3^{r-1}\sigma
\]
give \(\sigma\mid3^{r-1}y\).
The legal affine instruction decreases the first component to \(r-1\),
regardless of the increase in \(y\). Its cost is one affine instruction;
the induction hypothesis supplies at most \(r-1\) additional ones.
In both branches, prepend the legal step to the finite inductive prefix.
This proves entry, the affine budget and compatibility of the same path.

Initially take \(r=b\), since \(\sigma=3^b\mid3^b n\).
Every state is positive: legal positive even division stays positive,
and \(3x+\sigma>0\) on the affine branch. At the entered state
\(\sigma m>0\), the quotient \(m\) is therefore positive.
\end{proof}

% The entry witness is chosen once and retained for all future times.
\begin{theorem}[Common compressed and original tails]\label{CommonTail}
For every \(n>0\), one pair \(t\in\N,m>0\) satisfies
\[
G_b^t(n)=\sigma m,\qquad A_{G_b}(t,n)\le b,
\]
and, for every \(s\in\N\),
\[
G_b^{t+s}(n)=\sigma H^s(m).
\]
At the original entry time \(r=\tau_{a,b}(t,n)\), the same \(m\) satisfies
\[
F_{a,b}^{r+s}(n)=\sigma F_{a,0}^{s}(m)
\qquad\text{for every }s\in\N.
\]
\end{theorem}

\begin{proof}
Choose \(t,m\) once by Lemma~\ref{FiniteEntry}.
Iterate composition and Lemma~\ref{ScaleOdd} give
\(G_b^{t+s}(n)=G_b^s(\sigma m)=\sigma H^s(m)\) for every \(s\).
Lemma~\ref{BlockClock} gives \(F_{a,b}^{r}(n)=\sigma m\).
Applying the second identity of Lemma~\ref{ScaleOdd} at this same entered
state yields \(F_{a,b}^{r+s}(n)=\sigma F_{a,0}^{s}(m)\).
No subsequent choice of an entry witness is made.
\end{proof}

The second tail uses the original raw time \(s\).
Comparison with the classical map instead has the clock
\[
F_{a,b}^{\,r+u+aR_0(u,m)}(n)=\sigma H^u(m).
\]
Indeed apply Lemma~\ref{BlockClock} with \(b=0\) to the right-hand side
of the original tail. The two clock statements are distinct.
The entered seed \(m\) is not asserted to equal the original \(n\).

\section{The remaining positive capture obligation}

% Both implications retain the exact targets and their positive domains.
\begin{theorem}[Positive capture]\label{CaptureEquivalence}
For each fixed \(a,b\in\N\),
\[
\begin{aligned}
&\bigl(\forall n\in\Z,\ n>0\Longrightarrow
          \exists r\in\N:F_{a,b}^r(n)=3^b\bigr)\\
&\qquad\Longleftrightarrow
\bigl(\forall m\in\Z,\ m>0\Longrightarrow
          \exists u\in\N:H^u(m)=1\bigr).
\end{aligned}
\]
\end{theorem}

\begin{proof}
Assume first the family assertion. Given any \(m>0\), the seed
\(\sigma m\) is positive. Its family orbit hits \(\sigma\).
Theorem~\ref{TargetReflection} yields a compressed hit
\(G_b^k(\sigma m)=\sigma\). Lemma~\ref{ScaleOdd} identifies that state
with \(\sigma H^k(m)\). Since \(\sigma\ne0\), cancellation gives
\(H^k(m)=1\). This holds for every positive classical seed \(m\).

Conversely assume the classical assertion and take an arbitrary \(n>0\).
Use the common entry \(G_b^t(n)=\sigma m\) from
Theorem~\ref{CommonTail}, where \(m>0\).
There is \(u\in\N\) with \(H^u(m)=1\).
The compressed tail gives \(G_b^{t+u}(n)=\sigma\), and
Theorem~\ref{TargetReflection} recovers an original raw hit.
The seed \(n\) was arbitrary, so the family assertion follows.
\end{proof}

This theorem proves an equivalence of universal positive obligations.
It does not prove the classical assertion and does not assert a pointwise
comparison with classical capture at the same original seed.
The finite entry obligations are solved by Lemma~\ref{FiniteEntry};
the remaining universal obligation is precisely the one displayed
in Theorem~\ref{CaptureEquivalence}.

\section{The designated primitive original cycle}

% All intermediate division guards and every shorter positive return are checked.
\begin{theorem}[Primitive raw period]\label{PrimitiveCycle}
The state \(\sigma\) has primitive original raw period \(a+3\).
Its complete trace is
\[
\sigma,\ 2^{a+2}\sigma,\ 2^{a+1}\sigma,\ \ldots,\ 2\sigma,\ \sigma,
\]
and more precisely
\[
F_{a,b}^{j+1}(\sigma)=2^{a+2-j}\sigma
\qquad(0\le j\le a+2).
\]
\end{theorem}

\begin{proof}
Oddness of \(\sigma\) forces its affine branch, whose image is
\(2^a(3\sigma+\sigma)=2^{a+2}\sigma\).
Each subsequent positive binary exponent gives an even state and hence
a legal division by \(2\). Induction proves the trace through exponent zero,
which occurs at time \(a+3\).

If \(0<r<a+3\) were a return, put \(j=r-1\).
The trace would give
\(\sigma=2^{a+2-j}\sigma\), with \(a+2-j>0\).
Canceling \(\sigma\ne0\) forces a strictly positive power of \(2\) to
equal one, a contradiction. Thus every positive return time is at least
\(a+3\), proving primitive minimality against all original raw returns.
No universal convergence assumption enters this argument.
\end{proof}

\section{Concrete instances and finite replay}

The exponent-zero instance has the raw cycle \(1,4,2,1\).
At \((a,b)=(2,1)\), the designated cycle is
\[
3,\ 48,\ 24,\ 12,\ 6,\ 3.
\]
The seed \(1\) enters \(3\Z\) along \(1,24,12,6\), at original time three
with normalized seed two, and then reaches the target three.
Its compressed prefix is \(1,6\), with one affine instruction.

At \((a,b)=(1,2)\), the designated cycle is \(9,72,36,18,9\).
The original seed-one prefix is
\[
1,\ 24,\ 12,\ 6,\ 3,\ 36,\ 18.
\]
It enters \(9\Z\) at time six with normalized seed two.
The compressed prefix \(1,12,6,3,18\) uses two affine instructions,
and its expanded time is \(4+1\cdot2=6\).
The next original state is the target nine.

The exact signed replay uses arbitrary-precision integers and checks every
division guard, binary block endpoint, strict interior, entry affine budget,
same-clock original tail and separate classical variable clock.
Its main box is
\[
0\le a\le16,\qquad 0\le b\le12,\qquad1\le|n|\le64,
\]
containing \(17\cdot13\cdot128=28288\) seed certificates.
These are not distinct-orbit counts or an infinite capture proof.

The callable resource bounds are \(a\le64\), \(b\le32\) and at most
\(512\) seed bits. A compressed prefix cap of \(4096\), an original tail of
\(32\) steps and a classical comparison of \(16\) steps bound computation
only. A refused prefix is not a noncapture certificate.
Zero, nil, oversized and unconfigured seed inputs are refused.

Eight chunks have disjoint workspaces and result slots.
Operands reserve \(2048\) bits before replay.
For the allowed ternary exponent, \(3^{32}=1853020188851841<2^{51}\).
The entry budget permits at most \(32\) compressed affine instructions,
each satisfying
\[
|3x+\sigma|\le4\max\{|x|,\sigma\}.
\]
Thus compressed prefix states are smaller in magnitude than \(2^{576}\),
and binary interiors than \(2^{640}\).
After entry, allowing the scale separately gives the conservative bounds
\(2^{755}\) for the original \(32\)-step tail and \(2^{723}\) for the
expanded \(16\)-step classical comparison.
These operand bounds fit the reservation; they do not assert universal
zero allocation. Warmed zero allocation is measured only at the four
tested pairs \((0,0),(16,12),(64,32),(3,0)\) with seed \(27\).
Cycle replay covers all \(65\cdot33=2145\) bounded exponent presentations,
not all cycles.

\section{Formal correspondence and domain boundaries}

The formal file is \texttt{ClassicalAbsorption.lean}.
Its parameter definition retains signed integer coefficients and the
division-first guard. The principal correspondences are:
\begin{itemize}
\item \texttt{ClassicalAbsorptionCompressionCaptureIff}:
      Theorem~\ref{TargetReflection}, for every signed seed.
\item \texttt{ClassicalAbsorptionCompressedTail}:
      compressed entry, affine budget and one common classical tail.
\item \texttt{ClassicalAbsorptionRawTail}:
      the original raw tail, Theorem~\ref{CommonTail}.
\item \texttt{ClassicalAbsorptionPositiveCaptureIff}:
      positive capture, Theorem~\ref{CaptureEquivalence}.
\item \texttt{ClassicalAbsorptionCycleTrace}:
      the explicit division phases, Theorem~\ref{PrimitiveCycle}.
\item \texttt{ClassicalAbsorptionPrimitiveCycle}:
      minimal return time, Theorem~\ref{PrimitiveCycle}.
\end{itemize}
The finite-entry proof uses the existing guarded offset-entry induction;
the raw-tail proof uses the existing binary endpoint identity and
coprime scaling on the invariant lattice.
The formal first-recursion affine count used in the entry budget and
the endpoint-recursion compression count are separate APIs.
No additional formal equality of these APIs is asserted here.

Finally, the broader source family \((3c,c,2)\), with \(c=2^a u\) and odd
\(u\), compresses to \((3u,u,2)\). Coprime offset normalization leaves
\((3u,1,2)\), which is classical only when \(u=1\).
It is not covered by silently replacing that multiplier by \(3\).

\bibliographystyle{amsalpha}
\bibliography{ClassicalAbsorption}
\end{document}
